\begin{center}
    \large
    
    \begin{singlespace}
        \textbf{\englishTitle{}} \\[0.5cm]
    \end{singlespace}
    
    \begin{singlespace}
        Student : \studentEnName{}  \hspace{1.0cm} 
        % 兩個指導教授要寫 Advisors
        \ifdefined\advisorCnNameB
            Advisors: Dr.\, \advisorEnName \\
            \hspace{6.6cm} Dr.\, \advisorEnNameB  \\
        \else
            Advisor: Dr.\, \advisorEnName \\
        \fi
    \end{singlespace}
    
    \vspace{0.5cm}
    \begin{singlespace}
        \NameofDepartmentInstituteEN\\
        National Yang Ming Chiao Tung University\\[0.2cm]
    \end{singlespace}
    
    \textbf{Abstract} \\[0.5cm]

\end{center}

\normalsize 
  
In previous studies, researchers commonly used gradient-based adjoint methods (Adjoint Method) for the design of photonic structures.
However, this method easily falls into local optima, making it difficult to explore globally optimal solutions.
Therefore, this study aims to combine machine learning and quantum annealing techniques, and introduce encoding methods to search for globally optimal structures, providing a new approach for photonic structure design.
This study integrates Factorization Machines (FM) and Quantum Annealing (QA) to construct an FMQA framework and applies it to the inverse design of two-dimensional silicon photonic waveguide structures.

However, due to the limited number of qubits in current quantum annealing hardware, 
this study proposes multi-encoding methods to enhance the representational capability of
 design parameters and the resolution of the design space. This approach also mitigates 
 discretization-induced quantization errors and the staircase effect, 
 ensuring high consistency between numerical models and real physical behavior. 
 In terms of spatial and feature encoding, three encoding strategies are employed: 
 Gaussian Upsampling, two-dimensional Fourier integral, and beta-Variational Autoencoder ($\beta$-VAE).

For numerical simulations, the finite-difference time-domain (FDTD) method is used to analyze the 
propagation behavior of transverse electric (TE) multimodes in two-dimensional waveguide components. 
The figure of merit (FoM) is set as the optimization objective to constrain transmission loss 
and improve transmission efficiency and mode purity.

This study designs and analyzes two key silicon photonic components: 
a 90-degree Bend Waveguide and a Compact Waveguide Crossing. In the bend waveguide design, 
the primary objective is to reduce transmission loss and improve transmission efficiency. 
In the compact waveguide crossing, structural symmetry is exploited to reduce the number of 
required qubits, effectively suppress mode conversion and crosstalk while maintaining 
high transmission efficiency. Through the proposed FMQA framework and customized objective function, 
this study successfully derives high-performance optimized waveguide structures with support 
for multimode transmission, demonstrating new possibilities and directions for 
future silicon photonic integrated circuit design.


\vspace{1cm}

% 5-7 Keywords (English) 
\textbf{Keywords: Factorization Machine, Silicon Photonics, Waveguide, Inverse Design, Quantum Annealing, Machine Learning} 
